\documentclass[11pt]{article}
\usepackage[margin=25mm]{geometry}
\usepackage{fontspec}
\setmainfont{Latin Modern Roman}
\usepackage{amsmath,amssymb,mathtools}
\usepackage{unicode-math}
\setmathfont{Latin Modern Math}
\usepackage{xurl,hyperref}
\hypersetup{colorlinks=true,urlcolor=blue}
\setcounter{secnumdepth}{0}
\setlength{\parindent}{0pt}
\setlength{\parskip}{5pt}
\setlength{\emergencystretch}{3em}
\newcommand{\sref}[2]{\hyperlink{source:#1:#2}{[S#2]}}
\newcommand{\cataloguescope}[1]{\paragraph{Recorded target.}#1}
\newcommand{\R}{\mathbb R}
\newcommand{\C}{\mathbb C}
\newcommand{\N}{\mathbb N}
\newcommand{\Z}{\mathbb Z}
\newcommand{\Q}{\mathbb Q}
\newcommand{\F}{\mathbb F}
\newcommand{\eps}{\varepsilon}
\begin{document}
% UnsolvedMath ID 30004094; source catalogue code OWR-16772-001
\setcounter{section}{30004093}
\section{Higher-Order Failure of the Quantum Linear Boltzmann Limit}\label{30004094}
{\small UnsolvedMath ID 30004094 · Mathematical Physics}
\subsection{Definitions and mathematical statement}

\paragraph{Shared notation.}$\mathbb N=\{1,2,\ldots\}$; the standard number fields have their usual meanings. All additional conventions are stated locally.


\paragraph{Definitions and conventions.}
The quantum Lorentz gas describes a quantum particle scattering from fixed scatterers; here their centers lie periodically on $\mathbb Z^d$. The Boltzmann--Grad limit is a dilute-scattering scaling of the evolution or its phase-space distribution. A coupling-constant expansion sorts terms by the strength of an individual scatterer. Agreement at an order means equality of the corresponding coefficient with that obtained from the linear Boltzmann transport equation, using the same scaling and initial-data class.

\paragraph{Target retained from the canonical source catalogue.}
\paragraph{Catalogue wording retained for traceability.}
For the quantum Lorentz gas with scatterers placed periodically on $\mathbb Z^d$, does the Boltzmann–Grad limit agree with the linear Boltzmann equation through second order in the coupling constant but disagree at every order greater than four?

\paragraph{Source scope and qualification.}
The question singles out orders greater than four and does not state an answer for third or fourth order. Perturbative agreement through second order is distinct from convergence of the full evolution at fixed coupling.
\subsection{Short English statement}
Does periodic quantum scattering agree with Boltzmann transport to second order but fail at each perturbative order above four?
\subsection{Research attempt}
\paragraph{Model, date and investigated scope.}
OpenAI GPT-6 Astra (configured model \texttt{gpt-6-astra}), ultra reasoning; generated 11 October 2026. Calculated lattice coherent peaks and their average to isolate a possible recollision mechanism; no claimed perturbative coefficient comparison.

\paragraph{Sources and literature check.}
Read the supplied definition, catalogue wording, source qualifications and reference entries for 30004094. The historical source/context and inherited status are retained; no fresh literature or novelty certification is claimed. No external theorem is used without proof in this attempt; explicit model assumptions are distinguished from facts about the source model.

\paragraph{A coherent-lattice calculation.} For \(N\) equally spaced scatterer locations \(j=0,\ldots,N-1\) in one dimension, the elementary coherent amplitude factor at momentum transfer \(q\) is \(S_N(q)=\sum_j e^{iqj}\). Summing a geometric series gives
\[
 |S_N(q)|^2=\frac{\sin^2(Nq/2)}{\sin^2(q/2)}
\]
away from multiples of \(2\pi\), with value \(N^2\) at those multiples by continuity. Expanding the squared sum and integrating exponentials proves \((2\pi)^{-1}\int_{-\pi}^{\pi}|S_N(q)|^2dq=N\). Hence a momentum average can look incoherent while narrow coherent peaks are much larger pointwise.
\paragraph{Attempt and first obstruction.} This suggests looking for lattice recollision contributions which survive the Boltzmann--Grad averaging, but it does not compute any coefficient of the source quantum dynamics. In particular one must track the peak widths, scaling, propagators and initial-data averaging before deciding whether the coherence survives. A nonzero discrepancy at one order also does not imply discrepancy at every higher order: the formal difference \(\lambda^6\) is zero at orders seven and above. Cancellations must be excluded separately at each requested order.
\paragraph{Scope audit.} Neither agreement through second order nor failure at every order greater than four is proved here. The third and fourth orders are left unspecified as in the source. The finite one-dimensional structure factor is a concrete test of an incoherent-averaging shortcut, not an asserted answer for the source dimension, scaling or full finite-coupling evolution.

\paragraph{Verification and adversarial audit.}
Checked the geometric sum, Bragg-limit value, integral normalization and the quantifier distinction between one discrepant order and every later order. This is a GPT-6 Astra Ultra self-audit, not independent human verification. Prompt A drafting followed by prompt B scope/dependency/boundary self-audit. Checked the geometric sum, Bragg-limit value, integral normalization and the quantifier distinction between one discrepant order and every later order. The remaining gap is: Compute the source-scaled recollision integrals and prove their noncancellation at each order above four, alongside second-order agreement. Unresolved status and zero original-target completion are retained.

Coordinator checked the finite geometric-series structure factor and its averaged squared norm. Higher-order noncancellation is expressly not inferred. This additional check was model-assisted, not independent human verification.

\paragraph{Tools, process and contributions.}
Python 3 and local file-reading/writing tools for manuscript serialization and structural checks; mathematical calculations were symbolic reasoning, with no numerical experiment or theorem prover. Contribution prompt A was followed by prompt B. GPT-6 Astra with ultra reasoning drafted the research and performed its stated self-audit; any additional coordinator contribution is identified in the verification and audit. A GPT-6-based Codex coordinator assembled the manuscript and checked repository integration; its exact variant and reasoning configuration were not exposed. No human mathematical contribution, independent human proof check or formal proof certification is supplied. Complete derivations and the consulted source references are included here; private reasoning transcripts are not published.

\paragraph{Final status and first remaining gap.}
\textbf{UNRESOLVED}. Compute the source-scaled recollision integrals and prove their noncancellation at each order above four, alongside second-order agreement.

\section{Completion Estimate}
\noindent\textbf{COMPLETION: 0\%}
The proved calculation is an elementary auxiliary result or a test of a restricted strategy; no new portion of the full source target is established.

\subsection{Sources}
\begin{itemize}
\item[\hypertarget{source:30004094:1}{S1}] ULAM.AI and the UnsolvedMath Contributors, UnsolvedMath, record 30004094 (OWR-16772-001). Catalogue curation under CC BY 4.0; definitions and English exposition added locally.
\url{https://huggingface.co/datasets/ulamai/UnsolvedMath}.
\item[\hypertarget{source:30004094:2}{S2}] Mini-Workshop: Lorentz Gas Dynamics: particle systems and scaling limits (2019), Oberwolfach Reports, DOI 10.4171/owr/2019/10, pp. 617--661. Primary report accessed 2026-10-11; relevant PDF pages 25.
\url{https://doi.org/10.4171/owr/2019/10}.
\item[\hypertarget{source:30004094:3}{S3}] Quantum Lorentz gas limits (catalogue-listed context; not independently consulted in this preparation).
\url{https://arxiv.org/abs/2201.08229}.
\end{itemize}
\label{end:30004094}
\end{document}
